%
% ratchet.tex
%
% Z Specification for the vox OO / coupling / suppression ratchet tools
% Formal model of the shared fail-closed, base-commit-authoritative core
% (unified in PR #312) and the three per-tool comparison deltas.
%
% This specification DOCUMENTS THE CODE AS IT IS. Where the model and the
% original informal brief disagree, the model follows the code; the
% divergences are recorded in section "Findings".
%

\documentclass[a4paper,10pt,fleqn]{article}
\usepackage[margin=1in]{geometry}
\usepackage{fuzz}
\usepackage[colorlinks=true,linkcolor=blue,citecolor=blue,urlcolor=blue]{hyperref}

\begin{document}

\title{The vox Ratchet Tools: A Z Specification}
\author{Formal model of check-oo, check-coupling, check-suppressions}
\date{July 2026}
\hypersetup{
  pdftitle={The vox Ratchet Tools: A Z Specification},
  pdfauthor={Punt Labs},
  pdfsubject={Z Specification},
  pdfcreator={fuzz/probcli}
}
\maketitle

\tableofcontents
\newpage

\section{Introduction}

Vox enforces three code-quality ratchets, each a command-line gate run
under \texttt{make check} and in CI:

\begin{itemize}
\item \textbf{check-oo} (\texttt{tools/oo\_ratchet}): a \emph{must-improve}
  ratchet on per-file object-orientation metrics.
\item \textbf{check-coupling} (\texttt{tools/coupling}): a
  \emph{regression-only} ratchet on per-file coupling/cohesion metrics.
\item \textbf{check-suppressions} (\texttt{tools/suppression}): a
  \emph{regression-only} ratchet on a single scalar --- the total count of
  lint/type suppression comments plus \texttt{per-file-ignores}.
\end{itemize}

PR~\#312 unified all three onto one \emph{fail-closed,
base-commit-authoritative} core. They differ only in what they compare.
This document therefore models the shared core once --- base resolution,
the fail-closed transition set, the CI-write guard, and the never-loosen
update contract --- and then adds one operation schema per tool for the
part that genuinely differs.

The specification type-checks under \texttt{fuzz -t}. Conjectures
(written \texttt{\textbackslash vdash?}) record the load-bearing
properties; \texttt{fuzz} checks that they are well-typed. Two of them
(sections \ref{sec:findings}) record where the code diverges from the
informal brief; they are stated as \emph{existence} conjectures because
they document behaviour the brief did not predict.

\section{Basic Types}

The three carriers the model needs no internal structure for: repository
relative Python file paths, metric names, and git commit hashes.

\begin{zed}
[FILE, METRIC, COMMIT]
\end{zed}

\section{Free Types}

The Boolean type: \texttt{fuzz} has none built in, and \texttt{ZBOOL}
avoids the \texttt{BOOL}/\texttt{true}/\texttt{false} keyword clash that
would defeat ProB.

\begin{zed}
ZBOOL ::= ztrue | zfalse
\end{zed}

A metric's ordering direction. \texttt{dGeq} means larger is better
(\texttt{>=} threshold), \texttt{dLeq} smaller is better (\texttt{<=}),
and \texttt{dEq} means closer to a fixed target is better (\texttt{==}).

\begin{zed}
Direction ::= dGeq | dLeq | dEq
\end{zed}

The four outcomes a ratchet operation can produce. \texttt{pass} and
\texttt{bootstrapPass} exit zero; \texttt{fail} and \texttt{failClosedError}
exit non-zero. \texttt{bootstrapPass} is reserved for a genuine first
adoption (no baseline anywhere); \texttt{failClosedError} is a controlled
non-zero raised when git or the baseline blob is broken.

\begin{zed}
OUTCOME ::= pass | bootstrapPass | fail | failClosedError
\end{zed}

\begin{axdef}
exitCode : OUTCOME \fun \nat
\where
exitCode = \{ pass \mapsto 0, bootstrapPass \mapsto 0,
fail \mapsto 1, failClosedError \mapsto 1 \}
\end{axdef}

The fail-closed characterisation of the exit code is fixed by the
definition above: \texttt{exitCode}$~o = 0$ exactly when $o = pass$ or
$o = bootstrapPass$. The two non-zero outcomes are the two failure kinds.

\section{Metric Ordering}

Each metric carries a direction and an absolute target. These are the
\texttt{Thresholds} tables in the code (\texttt{oo\_ratchet/thresholds.py},
\texttt{coupling/thresholds.py}). The model leaves the tables abstract:
only the ordering they induce matters to the ratchet.

\begin{axdef}
metricDir : METRIC \fun Direction \\
metricTarget : METRIC \fun \num
\end{axdef}

Metric values are integers here. The code uses floats; \texttt{fuzz} has
no reals, and the ratchet depends only on the \emph{ordering} of values,
never on their exact magnitude, so $\num$ carries the model faithfully.

The absolute value, defined without unary minus so \texttt{fuzz} and ProB
stay happy.

\begin{axdef}
zabs : \num \fun \nat
\where
\forall n : \num @ (n \geq 0 \implies zabs~n = n) \\
\quad~ \land (n < 0 \implies zabs~n + n = 0)
\end{axdef}

The distance of a value from a metric's target (used only by the
\texttt{dEq} direction).

\begin{axdef}
dist : METRIC \cross \num \fun \nat
\where
\forall m : METRIC; v : \num @ dist(m, v) = zabs(v - metricTarget~m)
\end{axdef}

``At least as good as'': the ratchet's non-regression test.

\begin{axdef}
betterEq : METRIC \cross \num \cross \num \fun ZBOOL
\where
\forall m : METRIC; c, b : \num @ betterEq(m, c, b) = ztrue \\
\quad~ \iff (metricDir~m = dGeq \land c \geq b) \\
\quad~ \lor (metricDir~m = dLeq \land c \leq b) \\
\quad~ \lor (metricDir~m = dEq \land dist(m, c) \leq dist(m, b))
\end{axdef}

``Strictly better'': the must-improve test.

\begin{axdef}
strictBetter : METRIC \cross \num \cross \num \fun ZBOOL
\where
\forall m : METRIC; c, b : \num @ strictBetter(m, c, b) = ztrue \\
\quad~ \iff (metricDir~m = dGeq \land c > b) \\
\quad~ \lor (metricDir~m = dLeq \land c < b) \\
\quad~ \lor (metricDir~m = dEq \land dist(m, c) < dist(m, b))
\end{axdef}

``Meets its absolute threshold'': used to grade a brand-new file that has
no baseline to improve against.

\begin{axdef}
meets : METRIC \cross \num \fun ZBOOL
\where
\forall m : METRIC; v : \num @ meets(m, v) = ztrue \\
\quad~ \iff (metricDir~m = dGeq \land v \geq metricTarget~m) \\
\quad~ \lor (metricDir~m = dLeq \land v \leq metricTarget~m) \\
\quad~ \lor (metricDir~m = dEq \land v = metricTarget~m)
\end{axdef}

\section{Metric Maps}

A file's metrics are a partial map from metric name to value; a baseline
is a partial map from file to metric map.

\begin{zed}
MMAP == METRIC \pfun \num
\end{zed}

\begin{zed}
BLINE == FILE \pfun MMAP
\end{zed}

The metrics of \texttt{c} that are worse than \texttt{b} --- the regressed
set. A metric absent from \texttt{b} never regresses (a new metric is not
a regression).

\begin{axdef}
worse : MMAP \cross MMAP \fun \finset METRIC
\where
\forall c, b : MMAP @ worse(c, b) = \\
\quad~ \{ m : \dom c | m \in \dom b \land betterEq(m, c~m, b~m) = zfalse \}
\end{axdef}

The metrics of \texttt{c} that strictly improve on \texttt{b}.

\begin{axdef}
better : MMAP \cross MMAP \fun \finset METRIC
\where
\forall c, b : MMAP @ better(c, b) = \\
\quad~ \{ m : \dom c | m \in \dom b \land strictBetter(m, c~m, b~m) = ztrue \}
\end{axdef}

Whether every metric of a map clears its absolute threshold.

\begin{axdef}
allMeet : MMAP \fun ZBOOL
\where
\forall c : MMAP @ allMeet~c = ztrue \iff \\
\quad~ (\forall m : \dom c @ meets(m, c~m) = ztrue)
\end{axdef}

\section{The Shared Fail-Closed Core}

Every check resolves a comparison base, reads the baseline committed
\emph{at that base commit}, and only then compares. The resolution and
the fail-closed decisions that precede the comparison are identical across
all three tools (\texttt{ratchet.py} for check-oo, \texttt{ratchet.py} for
check-coupling, \texttt{baseline.py} for check-suppressions).

\subsection{The git environment}

\texttt{GitEnv} records what the resolution phase can observe: whether the
base resolves, whether a baseline blob exists at the base and whether it is
corrupt, the same for the \texttt{origin/main} tip, whether an in-tree
baseline file is present, and whether \texttt{-{}-require-base} was passed.
The invariants encode the obvious well-formedness: a blob can be corrupt
only if it exists, and can exist only if its ref resolves.

\begin{schema}{GitEnv}
baseResolves : ZBOOL \\
baseHasBlob : ZBOOL \\
baseCorrupt : ZBOOL \\
tipResolves : ZBOOL \\
tipHasBlob : ZBOOL \\
tipCorrupt : ZBOOL \\
inTreeExists : ZBOOL \\
requireBase : ZBOOL
\where
baseCorrupt = ztrue \implies baseHasBlob = ztrue \\
baseHasBlob = ztrue \implies baseResolves = ztrue \\
tipCorrupt = ztrue \implies tipHasBlob = ztrue \\
tipHasBlob = ztrue \implies tipResolves = ztrue
\end{schema}

\subsection{Base unresolvable}

When no base commit can be found (\texttt{\_no\_base}): under
\texttt{-{}-require-base} the gate fails; a genuine first adoption --- no
in-tree baseline to protect --- bootstraps; otherwise an in-tree baseline
with an unresolvable base means a stale or unfetched \texttt{origin/main},
and the gate fails loudly rather than pass on trust.

\begin{schema}{NoBase}
GitEnv \\
out! : OUTCOME
\where
baseResolves = zfalse \\
requireBase = ztrue \implies out! = fail \\
requireBase = zfalse \land inTreeExists = zfalse \implies \\
\quad~ out! = bootstrapPass \\
requireBase = zfalse \land inTreeExists = ztrue \implies out! = fail
\end{schema}

\subsection{Base resolves but carries no baseline blob}

When the base commit has no baseline (\texttt{\_absent\_base\_baseline}):
a genuine first adoption requires the \texttt{origin/main} tip to also
lack a baseline. A corrupt tip blob is a controlled error. A tip that
carries a baseline means the branch forked before adoption and would
launder a regression past the empty base, so it fails. An unresolvable
tip with an in-tree baseline present cannot confirm first adoption, so it
fails closed.

\begin{schema}{AbsentBase}
GitEnv \\
out! : OUTCOME
\where
baseResolves = ztrue \land baseHasBlob = zfalse \\
tipCorrupt = ztrue \implies out! = failClosedError \\
tipResolves = zfalse \land inTreeExists = ztrue \implies out! = fail \\
tipResolves = zfalse \land inTreeExists = zfalse \implies \\
\quad~ out! = bootstrapPass \\
tipResolves = ztrue \land tipCorrupt = zfalse \land tipHasBlob = ztrue \\
\quad~ \implies out! = fail \\
tipResolves = ztrue \land tipHasBlob = zfalse \implies \\
\quad~ out! = bootstrapPass
\end{schema}

\subsection{Base baseline blob corrupt}

A base baseline that resolves and exists but fails to parse (or, for
coupling and suppression, is non-dict / non-numeric / non-UTF8) raises
\texttt{GitError}, which the CLI turns into a controlled non-zero.

\begin{schema}{BaseError}
GitEnv \\
out! : OUTCOME
\where
baseResolves = ztrue \land baseHasBlob = ztrue \land baseCorrupt = ztrue \\
out! = failClosedError
\end{schema}

\subsection{Reaching the comparison}

The comparison runs only when the base resolves and carries a readable
baseline blob. \texttt{Proceed} names that precondition; each per-tool
check schema includes it.

\begin{schema}{Proceed}
GitEnv
\where
baseResolves = ztrue \land baseHasBlob = ztrue \land baseCorrupt = zfalse
\end{schema}

\subsection{The CI-write guard}

Every mutation (update, relax, reconcile, rebaseline) refuses to write
under \texttt{GITHUB\_ACTIONS} unless \texttt{-{}-allow-ci-write} is
passed. The guard is the same in all three tools.

\begin{schema}{CiGuard}
githubActions : ZBOOL \\
allowCiWrite : ZBOOL \\
blocked! : ZBOOL
\where
blocked! = ztrue \iff \\
\quad~ (githubActions = ztrue \land allowCiWrite = zfalse)
\end{schema}

\subsection{Fail-closed properties}

The fail-closed contract is that \texttt{bootstrapPass} arises \emph{only}
from a genuine absence of baseline. These property schemas restate the
guarantee as type-checked predicates over the transition schemas; each is
a redundant strengthening of the schema it includes, so it has exactly the
same models --- it documents the property without weakening it.

\begin{schema}{NoBaseFailClosed}
NoBase
\where
out! = bootstrapPass \implies inTreeExists = zfalse
\end{schema}

\begin{schema}{AbsentBaseFailClosed}
AbsentBase
\where
out! = bootstrapPass \implies \\
\quad~ (inTreeExists = zfalse \\
\quad~ \lor (tipResolves = ztrue \land tipHasBlob = zfalse))
\end{schema}

\section{check-oo: The Must-Improve Ratchet}

The OO ratchet is the strict one. On the comparison path it requires, of a
change that touches any file already tracked in the base baseline, that at
least one metric \emph{strictly} improve on such a file, and that no metric
regress --- unless a regression is waived by an audited relaxation.

\subsection{Inputs}

\texttt{OoCheck} carries the resolved base commit and the base-commit
baseline (\texttt{baselineAt}$~$\texttt{baseCommit}, the authoritative
comparison source --- \emph{never} the in-tree file); the in-tree baseline
(read only for the lock-in and waiver checks); the current scores; the
touched set (touched Python files that scored clean); the touched files
that failed to parse; and the waivable \texttt{(file, metric)} pairs
harvested from audit relaxations recorded since the base commit.

\begin{schema}{OoCheck}
Proceed \\
baseCommit : COMMIT \\
baselineAt : COMMIT \pfun BLINE \\
baseB : BLINE \\
inTree : BLINE \\
current : BLINE \\
touched : \finset FILE \\
brokenTouched : \finset FILE \\
waivable : \finset (FILE \cross METRIC) \\
wvd : \finset (FILE \cross METRIC) \\
rawRegressed : \finset (FILE \cross METRIC) \\
regressedP : \finset (FILE \cross METRIC) \\
waivedP : \finset (FILE \cross METRIC) \\
missingF : \finset FILE \\
lockFailF : \finset FILE \\
basePresent : \finset FILE \\
improvedExisting : ZBOOL \\
newPasses : ZBOOL \\
improvementOk : ZBOOL \\
out! : OUTCOME
\where
baseCommit \in \dom baselineAt \\
baseB = baselineAt~baseCommit \\
touched \subseteq \dom current \\
touched \cap brokenTouched = \emptyset \\
wvd = \{ f : touched; m : METRIC | f \in \dom inTree \land m \in \dom (inTree~f) \land inTree~f~m = current~f~m \land (f, m) \in waivable \} \\
rawRegressed = \{ f : touched; m : METRIC | f \in \dom baseB \land m \in worse(current~f, baseB~f) \} \\
regressedP = rawRegressed \setminus wvd \\
waivedP = rawRegressed \cap wvd \\
missingF = \{ f : touched | f \notin \dom inTree \} \\
lockFailF = \{ f : touched | f \in \dom inTree \land (\exists m : \dom (current~f) @ m \notin \dom (inTree~f) \lor inTree~f~m \neq current~f~m) \} \\
basePresent = touched \cap \dom baseB \\
improvedExisting = ztrue \iff \\
\quad~ (\exists f : basePresent @ better(current~f, baseB~f) \neq \emptyset) \\
newPasses = ztrue \iff \\
\quad~ (\exists f : touched \setminus \dom baseB @ allMeet(current~f) = ztrue) \\
improvementOk = ztrue \iff \\
\quad~ waivedP \neq \emptyset \\
\quad~ \lor (basePresent \neq \emptyset \land improvedExisting = ztrue) \\
\quad~ \lor (basePresent = \emptyset \land newPasses = ztrue) \\
brokenTouched \neq \emptyset \implies out! = fail \\
brokenTouched = \emptyset \land touched = \emptyset \implies out! = pass \\
brokenTouched = \emptyset \land touched \neq \emptyset \implies \\
\quad~ (out! = pass \iff \\
\quad~ (missingF = \emptyset \land lockFailF = \emptyset \\
\quad~ \land regressedP = \emptyset \land improvementOk = ztrue)) \\
out! \neq bootstrapPass \land out! \neq failClosedError
\end{schema}

The waiver is deliberately narrow: a metric is waived (\texttt{RELAXED})
only when the in-tree baseline has locked the current value (so the
regression was recorded, not smuggled) \emph{and} the audit log carries a
matching relaxation. The improvement gate follows the code's S5: if any
touched file is tracked in the base, the required improvement must come
from such a file; a pure-add change is satisfied by a new file clearing all
absolute thresholds; a change carrying a waiver is exempt.

\section{check-coupling: The Regression-Only Ratchet}

The coupling ratchet drops the must-improve gate, the waiver, the lock-in
check, and the completeness check. A touched file fails only if a metric is
strictly worse than its base-commit baseline; holding steady passes.

\subsection{The empty-baseline intercept}

Coupling adds one guard the OO ratchet lacks. An empty \texttt{\{\}} base
baseline (a truncated write or a bad merge) under \texttt{-{}-require-base}
fails closed before the comparison: otherwise every touched file would look
new and no regression could be seen.

\begin{schema}{CplEmptyRequire}
Proceed \\
baseCommit : COMMIT \\
baselineAt : COMMIT \pfun BLINE \\
baseB : BLINE \\
out! : OUTCOME
\where
baseCommit \in \dom baselineAt \\
baseB = baselineAt~baseCommit \\
baseB = \emptyset \\
requireBase = ztrue \\
out! = fail
\end{schema}

\subsection{The comparison}

The comparison runs when the base baseline is non-empty, or empty without
\texttt{-{}-require-base} (in which case every file is new and passes). A
file absent from the base is new: it never regresses.

\begin{schema}{CplCheck}
Proceed \\
baseCommit : COMMIT \\
baselineAt : COMMIT \pfun BLINE \\
baseB : BLINE \\
current : BLINE \\
touched : \finset FILE \\
brokenTouched : \finset FILE \\
regressedP : \finset (FILE \cross METRIC) \\
out! : OUTCOME
\where
baseCommit \in \dom baselineAt \\
baseB = baselineAt~baseCommit \\
\lnot (baseB = \emptyset \land requireBase = ztrue) \\
touched \subseteq \dom current \\
touched \cap brokenTouched = \emptyset \\
regressedP = \{ f : touched; m : METRIC | f \in \dom baseB \land m \in worse(current~f, baseB~f) \} \\
brokenTouched \neq \emptyset \implies out! = fail \\
brokenTouched = \emptyset \land touched = \emptyset \implies out! = pass \\
brokenTouched = \emptyset \land touched \neq \emptyset \implies \\
\quad~ (out! = pass \iff regressedP = \emptyset) \\
out! \neq bootstrapPass \land out! \neq failClosedError
\end{schema}

\section{check-suppressions: The Scalar-Count Ratchet}

The suppression ratchet compares a single scalar: the total count of
\texttt{noqa}, \texttt{type: ignore}, \texttt{pylint: disable},
\texttt{pyright: ignore} comments on code lines, plus the number of ruff
\texttt{per-file-ignores} rule codes. The check fails only if the current
total exceeds the base-commit total. The base total is coerced defensively
(\texttt{\_as\_int}): a non-numeric, boolean, \texttt{NaN}, or infinite
value becomes $0$, which fails closed --- a smaller base makes the current
total register as an increase.

\begin{schema}{SupCheck}
Proceed \\
baseCommit : COMMIT \\
supBaselineAt : COMMIT \pfun \nat \\
baseTotal : \nat \\
curTotal : \nat \\
out! : OUTCOME
\where
baseCommit \in \dom supBaselineAt \\
baseTotal = supBaselineAt~baseCommit \\
curTotal > baseTotal \implies out! = fail \\
curTotal \leq baseTotal \implies out! = pass \\
out! \neq bootstrapPass \land out! \neq failClosedError
\end{schema}

\section{The Never-Loosen Update}

\texttt{-{}-update} tightens a baseline but never loosens it. The base ref
scopes the touched set only; the regression test is against the
\emph{in-tree} baseline. check-oo and check-coupling share the algorithm
(\texttt{apply.py}, \texttt{writer.py}); check-suppressions does not (see
Findings).

The rename-aware base lookup: a file inherits its predecessor's entry when
it is a rename target absent from the baseline under its new name, so a
regression cannot launder itself through a rename. A file with no entry and
no rename predecessor is new; the empty metric map models ``no entry''
(its \texttt{worse} set is empty, so a new file never regresses).

\begin{axdef}
baseFor : BLINE \cross (FILE \pinj FILE) \cross FILE \fun MMAP
\where
\forall b : BLINE; r : FILE \pinj FILE; f : FILE @ \\
\quad~ (f \in \dom b \implies baseFor(b, r, f) = b~f) \\
\quad~ \land (f \notin \dom b \land f \in \dom r \land r~f \in \dom b \implies baseFor(b, r, f) = b~(r~f)) \\
\quad~ \land (f \notin \dom b \land \lnot (f \in \dom r \land r~f \in \dom b) \implies baseFor(b, r, f) = \emptyset)
\end{axdef}

\subsection{The update operation}

The in-tree baseline is the state. A touched file is \emph{writable} when
it scored clean and does not regress against its rename-aware in-tree base;
\emph{refused} when it scored clean but regresses (its old entry is kept,
so the carried base survives to the next run); dropped when it is gone from
disk. Writing a rename target drops the stale source key. The write refuses
under the CI guard, and fails the operation if anything was refused or
unparseable.

\begin{schema}{NeverLoosenUpdate}
inTree : BLINE \\
inTree' : BLINE \\
current : BLINE \\
touched : \finset FILE \\
renames : FILE \pinj FILE \\
parseErrTouched : \finset FILE \\
githubActions : ZBOOL \\
allowCiWrite : ZBOOL \\
guarded : ZBOOL \\
writable : \finset FILE \\
refusedF : \finset FILE \\
droppedMiss : \finset FILE \\
droppedSet : \finset FILE \\
updates : BLINE \\
out! : OUTCOME
\where
guarded = ztrue \iff (githubActions = ztrue \land allowCiWrite = zfalse) \\
writable = \{ f : touched | f \notin parseErrTouched \land f \in \dom current \land worse(current~f, baseFor(inTree, renames, f)) = \emptyset \} \\
refusedF = \{ f : touched | f \notin parseErrTouched \land f \in \dom current \land worse(current~f, baseFor(inTree, renames, f)) \neq \emptyset \} \\
droppedMiss = \{ f : touched | f \notin parseErrTouched \land f \notin \dom current \} \\
droppedSet = droppedMiss \cup \ran (writable \dres renames) \\
updates = writable \dres current \\
guarded = ztrue \implies (out! = fail \land inTree' = inTree) \\
guarded = zfalse \implies \\
\quad~ inTree' = (droppedSet \ndres inTree) \oplus updates \\
guarded = zfalse \land refusedF = \emptyset \land parseErrTouched = \emptyset \\
\quad~ \implies out! = pass \\
guarded = zfalse \land (refusedF \neq \emptyset \lor parseErrTouched \neq \emptyset) \\
\quad~ \implies out! = fail \\
out! \neq bootstrapPass \land out! \neq failClosedError
\end{schema}

\subsection{The never-loosen property}

The guarantee: no metric of any file present both before and after the
update is worse afterwards. This property schema restates it as a
type-checked predicate over the operation.

\begin{schema}{NeverLoosens}
NeverLoosenUpdate
\where
\forall f : \dom inTree' \cap \dom inTree @ \\
\quad~ worse(inTree'~f, inTree~f) = \emptyset
\end{schema}

\section{The Suppression Update: A Divergence}
\label{sec:findings}

check-suppressions' \texttt{update} does \emph{not} implement never-loosen.
After the CI guard it writes the current total unconditionally --- it takes
no base ref, does no scoping, and refuses no increase. The written total
may exceed the prior total.

\begin{schema}{SupUpdate}
supTotal : \nat \\
supTotal' : \nat \\
curTotal : \nat \\
githubActions : ZBOOL \\
allowCiWrite : ZBOOL \\
guarded : ZBOOL \\
out! : OUTCOME
\where
guarded = ztrue \iff (githubActions = ztrue \land allowCiWrite = zfalse) \\
guarded = ztrue \implies (out! = fail \land supTotal' = supTotal) \\
guarded = zfalse \implies (out! = pass \land supTotal' = curTotal)
\end{schema}

\section{Findings}

Modelling the code faithfully surfaced three places where behaviour differs
from the informal brief. The model follows the code.

\subsection{The OO ratchet does not harden its base baseline}

The brief lumped all three tools together: ``corrupt / non-dict / non-UTF8
/ NaN / bool baseline $\rightarrow$ a controlled fail; none may yield
pass.'' That holds for check-coupling and check-suppressions, whose blob
readers reject a non-dict, non-numeric, boolean, or non-UTF8 baseline as a
\texttt{GitError} (a controlled non-zero). It does \emph{not} hold for
check-oo: \texttt{oo\_ratchet/gitio.py}'s \texttt{show\_baseline} catches
only \texttt{json.JSONDecodeError}. A parseable-but-ill-typed base baseline
(a string value where a metric map is expected) is not rejected; the
comparison then indexes a string and raises an \emph{uncaught}
\texttt{TypeError} --- not the controlled \texttt{fail} the brief assumes.
In the model, \texttt{BaseError} covers only the unparseable case for
check-oo; the ill-typed-but-parseable case has no controlled outcome.

\subsection{The OO ratchet has no empty-baseline intercept}

check-coupling fails closed on an empty \texttt{\{\}} base baseline under
\texttt{-{}-require-base} (\texttt{CplEmptyRequire}). check-oo has no such
guard: an empty base baseline flows straight into the comparison, where
every touched file is treated as new. A pure-add change whose new files all
clear their absolute thresholds then \emph{passes} --- not as a genuine
first-adoption \texttt{bootstrapPass}, but as an ordinary \texttt{pass}.
\texttt{OoCheck} admits exactly this: with $baseB = \emptyset$,
$basePresent = \emptyset$, and \texttt{newPasses} = ztrue, the verdict is
\texttt{pass}. check-suppressions, by contrast, fails on an empty base
because the coerced base total is $0$ and any current suppression exceeds
it.

\subsection{Suppression update is not never-loosen}

As above (\texttt{SupUpdate}): check-oo and check-coupling refuse any
per-metric regression on update; check-suppressions rewrites the total
unconditionally once past the CI guard. \texttt{SupUpdate} permits
$supTotal' > supTotal$ with $out! = pass$.

\section{System Invariants}

The three properties the unified core guarantees, and where each lives in
this model:

\begin{itemize}
\item \textbf{Fail-closed} --- \texttt{bootstrapPass} arises only from a
  genuine absence of baseline: \texttt{NoBaseFailClosed},
  \texttt{AbsentBaseFailClosed}, and the \texttt{exitCode} definition.
\item \textbf{Never-loosen} --- an update's written baseline is no worse on
  any metric: \texttt{NeverLoosens} (for check-oo and check-coupling; see
  Findings for check-suppressions).
\item \textbf{Base-commit-authoritative} --- the comparison reads
  $\mathit{baselineAt}~\mathit{baseCommit}$, never the in-tree baseline:
  the equality $baseB = baselineAt~baseCommit$ in \texttt{OoCheck},
  \texttt{CplCheck}, and (as a total) \texttt{SupCheck}. The in-tree
  baseline appears in \texttt{OoCheck} only for the lock-in and waiver
  tests, and in \texttt{NeverLoosenUpdate} as the never-loosen target.
\end{itemize}

\end{document}
